\paragraph{AQ-136-FREE-ZERO-FORMS: a weight-one retraction still obstructs a full free PD de Rham source.}
The earlier AQ-136-DERHAM-CELLS excluded free zero-form
presentations whose new variables all have weight at least two.
The following refinement permits arbitrary variables of weight
one. Its additional hypothesis is automatic for a map to the
actual admissible truncated standard theory described below.

Let $k$ be a commutative ring, $S=\Gamma_k(T)$, and let $P_*$
be a simplicial graded PD $S$-algebra with homogeneous
presentations
\[
                  P_n=S\langle X_n\rangle_{\rm PD},
                \qquad \operatorname{wt}(x)\ge1.
                                                               \tag{136.FK1}
\]
Simplicial operators are arbitrary homogeneous PD $S$-maps.
Suppose its weight-one simplicial module has a coefficient
retraction
\[
          \rho:P_*^1\longrightarrow\underline{kT},
                            \qquad \rho(T)=T.              \tag{136.FK2}
\]
For every prime $p$ with $k/pk\ne0$, one has
\[
 [T^{[p]}]=0\text{ in }H_0(NP_*^p)
       \quad\Longrightarrow\quad
 H_1\!\left(N\Omega^1_{P_*/S,\mathrm{sp}}{}^{\,p}\right)\ne0.
                                                               \tag{136.FK3}
\]
Thus, if infinitely many rational primes are nonunits in $k$,
the full relative special-PD de Rham algebra of such a $P_*$
cannot be admissible. No flatness of Frobenius or of $k/pk$
is assumed.

\begin{proof}
Put $R=k/pk$ and $\bar P=P\otimes_kR$. As in the earlier
proof, form the weight-$p$ ordinary indecomposables and the
augmented special differentials
\[
 Q_*=(\bar P_+/(\bar P_+)^2)^p,\qquad
 F_*=\left(R\otimes_{\bar P_*}
          \Omega^1_{\bar P_*/\Gamma_R(T),\mathrm{sp}}\right)^p.
                                                               \tag{136.FK4}
\]
The augmentation sends all positive weights to zero.
These are simplicial $R$-modules. The special derivation,
followed by augmentation, induces a natural map $Q_*\to F_*$.

In every simplex the PD monomial basis gives the following
basis for $Q_n$:
\[
 T^{[p]},\qquad x^{[p]}\quad(\operatorname{wt}(x)=1),
       \qquad y\quad(\operatorname{wt}(y)=p).               \tag{136.FK5}
\]
Indeed a monomial involving at least two distinct base
variables is a product of two positive-weight monomials
with coefficient one. A pure divided power of a variable
of weight $a$ can have weight $p$ only when $(a,m)=(1,p)$
or $(a,m)=(p,1)$. Products of positive pure powers with total
exponent $p$ have the coefficient $\binom pi=0$ in $R$.
Consequently none of the pure vectors in (136.FK5) is killed
or related by the ordinary product ideal. This argument
works over an arbitrary ring $R$ of characteristic $p$.

The module $F_n$ is free on the $dy$ with
$\operatorname{wt}(y)=p$. The vector $y$ in (136.FK5) maps
to $dy$, while all the other displayed vectors map to zero:
$dT^{[p]}=0$ relatively, and
\[
                  d(x^{[p]})=x^{[p-1]}dx
\]
vanishes after the coefficient augmentation. Thus the kernel
$K_n$ is free on $T^{[p]}$ and the $x^{[p]}$ of weight-one
variables, and there is an exact sequence
\[
                  0\longrightarrow K_*
                      \longrightarrow Q_*
                      \longrightarrow F_*
                      \longrightarrow0.                  \tag{136.FK6}
\]
The point requiring care is the simplicial structure on
$K_*$ when faces mix the displayed weight-one variables.

Let $W_*=\bar P_*^1$ and let
$\mathrm{Fr}:R\to R$ denote $r\mapsto r^p$. Define
\[
              \mathrm{Fr}^*W_*=R\otimes_{\mathrm{Fr},R}W_*,
                                                               \tag{136.FK7}
\]
where the relation is
$a r^p\otimes w=a\otimes rw$.
The map
\[
                  w\longmapsto[\gamma_p(w)]\in Q_*
                                                               \tag{136.FK8}
\]
is additive and Frobenius-semilinear. Additivity follows
from the divided-power sum formula, since each mixed term
is a product of positive-weight elements. Semilinearity is
the identity $\gamma_p(rw)=r^p\gamma_p(w)$.
Consequently (136.FK8) induces a natural simplicial
$R$-linear map
\[
                    \mathrm{Fr}^*W_*\longrightarrow K_*.
                                                               \tag{136.FK9}
\]
The weight-one module $W_n$ is free on $T$ and the
weight-one variables. Formula (136.FK5) shows directly
that (136.FK9) takes a free basis to a free basis, so it is
an isomorphism. No preservation of exact sequences by
$\mathrm{Fr}^*$ is being used.

Reduce (136.FK2) modulo $p$ and apply this tensor functor
to the split pair
\[
          \underline{RT}\longrightarrow W_*
                                \xrightarrow{\bar\rho}\underline{RT}.
                                                               \tag{136.FK10}
\]
Every functor preserves this retraction identity. Also
$R\otimes_{\mathrm{Fr},R}R\simeq R$, by
$a\otimes b\mapsto ab^p$, with inverse $a\mapsto a\otimes1$.
Under (136.FK9), the resulting retraction of $K_*$
is a simplicial map to the constant module $R$ which sends
\[
                           z=T^{[p]}\longmapsto1.
                                                               \tag{136.FK11}
\]
In particular $[z]\ne0$ in $H_0(NK_*)$, because $R\ne0$.

The premise of (136.FK3) gives a normalized chain
$u\in N_1P_*^p$ with $\partial u=T^{[p]}$.
Its reduction and image in $Q_*$ is a normalized chain
$v$ with $\partial v=z$. Its image $a$ in $F_*$
is a cycle. The connecting homomorphism in (136.FK6) sends
$[a]$ to the nonzero class $[z]$, so $[a]\ne0$.
The integral relative differential $du$ is a normalized
one-cycle, since $\partial(du)=dT^{[p]}=0$.
Its reduction and augmented image is exactly $a$.
If $du$ were an integral boundary, $a$ would be a boundary,
a contradiction. This proves (136.FK3) without a flat
homology base-change assertion.

If the zero-form row is acyclic in a weight tail, its
coefficient class is killed at every sufficiently large
prime weight. Infinitely many nonunit primes therefore
give arbitrarily large nonzero one-form $H_1$ groups by
(136.FK3), contrary to horizontal admissibility.
\end{proof}

\paragraph{Actual-target consequence.}
Let $B$ be the standard admissible integral theory and fix
$N\ge2$. Define its strict PD subtheory
\[
 C_N^{w,*}=
 \begin{cases}
  kT^{[w]}[0],&w<N,\\
  B^{w,*},&w\ge N.
 \end{cases}                                               \tag{136.FK12}
\]
The coefficients in low weights are the existing coefficient
submodules of $B$. Products, differential and divided powers
preserve (136.FK12): a product involving a retained high
weight stays high, and positive divided powers multiply
weight by a positive integer. The vertical comparison is
unchanged in every sufficiently large weight, as is each
horizontal row. Thus $C_N\in\mathcal A_k$, exactly as in
AQ-136-WEIGHT-ONE.

Any strict map from the full relative special-PD de Rham
algebra of a $P_*$ satisfying (136.FK1) to $C_N$ induces
on weight-one zero forms the retraction (136.FK2), because
$C_N^{1,0}=\underline{kT}$ and coefficients are preserved.
If $k$ has infinitely many nonunit rational primes, the
source therefore cannot itself be admissible.
In particular, over $\mathbb Z$ even arbitrary additional
weight-one free zero-form cells cannot turn this full
relative-de-Rham construction into a source for $C_N$.

This is a construction-class obstruction. The target
$C_N$ is itself admissible and maps to itself, so it is
not a counterexample to the original common-source question.
The coskeleton replacements of AQ-136-FIRST-SKELETON need
not be degreewise free full de Rham algebras, and are not
excluded. Nonfree form relations and other strict PD
presentations remain available for higher comparisons.
